\chapter*{第三册能力清单}
\addcontentsline{toc}{chapter}{第三册能力清单}

第三册的出口是能把本地大模型推理拆成模型资产、张量、算子、KV cache、量化、compiler/runtime、后端、serving 和应用编排合同，并能用 oracle、trace 和 benchmark 证明关键判断。

\section*{能实现}

\begin{itemize}
  \item 实现 tiny decoder reference，覆盖 embedding、RMSNorm、Q/K/V、RoPE、GQA attention、KV append/read、SwiGLU、lm head 和 sampler。
  \item 实现模型导入闭环：tokenizer fixture、GGUF/safetensors 或等价 manifest、tensor name mapping、shape/dtype/quant verifier 和首 token trace。
  \item 实现 int8/int4 至少一种量化路径的 descriptor、packing、operator oracle、误差报告和容量报告。
  \item 实现 serving probe：admission、KV budget、prefill/decode 区分、TTFT、TPOT、p95、取消、backpressure 和拒绝原因。
\end{itemize}

\section*{能诊断}

\begin{itemize}
  \item 诊断 tokenizer/chat template/vocab/embedding 不一致导致的 token drift 或 logits drift。
  \item 诊断 KV cache 的 layer、position、kv\_head、dim、physical slot、prefix 和 sliding window 错误。
  \item 诊断 decode bytes/token、GQA/MQA 节省量、context length 斜率、paged KV 和 continuous batching 的资源影响。
  \item 诊断 AI compiler 中 shape inference、layout propagation、fusion legality、alias/liveness、buffer reuse 和 stateful op 调度约束。
\end{itemize}

\section*{能解释}

\begin{itemize}
  \item 解释 Transformer 数据流中 RMSNorm、RoPE、softmax、SwiGLU、sampling、低精度误差和数值稳定性的边界。
  \item 解释 GEMM/GEMV、prefill/decode、packing、blocking、register tile、cache tile、SIMD tail 和 GPU launch/copy/sync 的不同成本。
  \item 解释 llama.cpp/ggml、oneDNN、ONNX Runtime、vLLM、TensorRT-LLM、TVM/MLIR/XLA 等成熟系统应从哪些源码对象和设计合同拆解。
  \item 解释 LangChain/LangGraph/LlamaIndex 类上层框架如何影响 token budget、RAG trace、tool side effect、checkpoint 和本地推理服务接口。
\end{itemize}

\section*{不能夸大的部分}

第三册主线是本地推理引擎和 AI Infra 骨架。多 GPU 训练级别通信、生产级 CUDA/HIP kernel 库、完整分布式 serving 平台和所有模型格式方言仍需要专项项目继续扩展。
